\documentclass[12pt,a4paper]{article}

\usepackage{geometry}
\geometry{margin=1.2in}

\usepackage{amsmath}
\usepackage{booktabs}
\usepackage{hyperref}
\usepackage{microtype}
\usepackage{graphicx}
\usepackage{longtable}
\usepackage{array}
\usepackage{xcolor}

\title{Preliminary Analysis of Developer Activity and Language Ecosystems\\
in a Large GitHub Follow Network}
\author{}
\date{\today}

\begin{document}

\maketitle

\begin{abstract}
This report presents preliminary findings from a large-scale observational study
of developer activity within a GitHub follow network. Using the GitHub GraphQL
API, we collect activity metadata for accounts followed by a single user whose
follow graph exceeds one million accounts. Each account is classified as active
or inactive based on the recency of commits to public repositories, and the
dominant programming language of the most recently updated repository is
recorded.

Analysis of the first $\sim30{,}000$ accounts reveals systematic differences in
activity levels across programming language ecosystems, as well as evidence for
a characteristic ``half-life'' of developer activity on GitHub of approximately
ten years. These results provide an empirical snapshot of the structure and
longevity of open source developer communities.
\end{abstract}

\section{Introduction}

GitHub has become one of the central platforms for collaborative software
development, and its public data provides a unique opportunity to study the
dynamics of programming language ecosystems and developer activity over time.

Most existing studies rely on repository-level datasets such as GitHub Archive
or BigQuery mirrors of repository metadata. While these datasets capture
repository events in detail, they provide only indirect insight into the life
cycles of individual developers.

The present study instead analyzes a \emph{developer-centered dataset}. Starting
from a single GitHub account that follows more than one million other accounts,
we crawl the follow network using the GitHub GraphQL API and collect basic
metadata describing each followed account. The goal is to estimate:
\begin{itemize}
\item the fraction of developers who remain active over time,
\item differences in activity rates across programming languages,
\item and the approximate lifetime of developer participation on GitHub.
\end{itemize}

The results reported here are preliminary and based on the first $\sim30{,}000$
accounts processed during the crawl.

\section{Dataset Construction}

The dataset is collected using the GitHub GraphQL API. For each followed account
we retrieve:
\begin{itemize}
\item account creation timestamp,
\item number of followers,
\item number of public repositories,
\item the most recently pushed public repository,
\item and the primary language of that repository.
\end{itemize}

Accounts are classified as \textbf{active} if the most recent repository push
occurred within the previous year; otherwise the account is classified as
\textbf{inactive}. Accounts for which no language-tagged repository is available
appear in the dataset under the label \textsc{None}.

For the first $N=30{,}182$ accounts processed, the resulting dataset contains:

\begin{center}
\begin{tabular}{lr}
\toprule
Active accounts & 18,007 \\
Inactive accounts & 12,175 \\
\midrule
Active fraction & $\approx 0.60$ \\
\bottomrule
\end{tabular}
\end{center}

Because the sample is drawn from a follow network rather than uniformly across
GitHub, it is expected to be biased toward more active developers.

\section{Language Activity Ratios}

For each programming language we compute an \emph{activity ratio}
\[
R \;=\; \frac{\text{active}}{\text{active}+\text{inactive}}.
\]
This metric estimates the probability that a developer associated with a given
language ecosystem is still actively committing code. Table~\ref{tab:activity}
shows the complete results for all languages with at least 59 tagged accounts in
the dataset.

\begin{longtable}{lrrrc}
\toprule
Language & Active & Inactive & Total & Activity Ratio \\
\midrule
\endhead
\midrule
\multicolumn{5}{r}{\small\itshape continued on next page}\\
\endfoot
\bottomrule
\endlastfoot
\textsc{None}      & 3041 & 5155 & 8196 & 0.37 \\
Astro              &   54 &    8 &   62 & 0.87 \\
Rust               &  423 &   67 &  490 & 0.86 \\
TypeScript         & 2480 &  450 & 2930 & 0.85 \\
Dart               &  151 &   38 &  189 & 0.80 \\
Lua                &  114 &   26 &  140 & 0.81 \\
TeX                &   65 &   21 &   86 & 0.76 \\
Python             & 3366 & 1259 & 4625 & 0.73 \\
SCSS               &   92 &   38 &  130 & 0.71 \\
HTML               & 1386 &  599 & 1985 & 0.70 \\
Go                 &  484 &  203 &  687 & 0.70 \\
PowerShell         &   43 &   18 &   61 & 0.70 \\
R                  &  133 &   62 &  195 & 0.68 \\
Shell              &  415 &  208 &  623 & 0.67 \\
Kotlin             &  105 &   54 &  159 & 0.66 \\
Swift              &  125 &   66 &  191 & 0.65 \\
JavaScript         & 1695 &  954 & 2649 & 0.64 \\
Jupyter Notebook   &  695 &  383 & 1078 & 0.64 \\
C\#                &  277 &  192 &  469 & 0.59 \\
C++                &  476 &  346 &  822 & 0.58 \\
MATLAB             &   43 &   32 &   75 & 0.57 \\
Dockerfile         &   34 &   27 &   61 & 0.56 \\
CSS                &  228 &  200 &  428 & 0.53 \\
Vue                &   92 &   83 &  175 & 0.53 \\
PHP                &  188 &  174 &  362 & 0.52 \\
Java               &  507 &  527 & 1034 & 0.49 \\
Ruby               &  111 &  120 &  231 & 0.48 \\
Vim Script         &   27 &   33 &   60 & 0.45 \\
C                  &  307 &  400 &  707 & 0.43 \\
Objective-C        &    9 &   50 &   59 & 0.15 \\
\caption{Activity ratios for all languages with at least 59 accounts.
Languages are ordered by activity ratio descending.}
\label{tab:activity}
\end{longtable}

Several patterns are immediately visible.

Newer language ecosystems such as Rust, TypeScript, and Astro exhibit the
highest activity ratios, indicating strong ongoing development momentum. Astro
(0.87), in particular, reflects the recency of its adoption: the framework was
released publicly in 2021 and its community consists almost entirely of
developers who joined or began using GitHub recently. Rust (0.86) and TypeScript
(0.85) similarly attract developers who are active in contemporary engineering
contexts.

At the other end of the spectrum, older ecosystems show markedly lower
retention. Java (0.49) and Ruby (0.48) both fall below the overall mean,
reflecting the large populations of historical accounts accumulated since their
peaks of adoption. The \textsc{None} category---accounts for which no
language-tagged repository is available---shows the lowest ratio among
interpretable entries (0.37), consistent with the expectation that inactive
accounts are less likely to have recent, language-tagged repositories.

Objective-C (0.15) is an outlier in the full distribution and warrants separate
attention; it is discussed in Section~\ref{sec:objc}.

\subsection{The Objective-C Signal}\label{sec:objc}

The extremely low activity ratio for Objective-C ($R = 0.15$, 9 active out of
59 total) stands apart from all other ecosystems in the dataset. This figure is
consistent with the history of Apple platform development: following the
introduction of Swift at WWDC 2014 and its open-sourcing in 2015, the iOS and
macOS developer communities migrated wholesale to the new language over the
subsequent three to five years. Developers whose most recent language-tagged
repository is still Objective-C are, by definition, those who have not yet
updated to Swift---a population skewed heavily toward inactivity. The ratio thus
functions here not merely as a measure of ecosystem vitality but as a fossil
record of a completed platform migration.

\section{Account Age and Developer Survival}

Using account creation timestamps and most recent commit timestamps, we estimate
the fraction of developers remaining active as a function of account age.

\begin{center}
\begin{tabular}{cc}
\toprule
Account Age (years) & Active Fraction \\
\midrule
0  & 0.662 \\
5  & 0.670 \\
10 & 0.500 \\
11 & 0.478 \\
\bottomrule
\end{tabular}
\end{center}

The curve crosses $0.5$ at approximately ten years, suggesting a
\emph{developer activity half-life} of roughly a decade. This implies that half
of the developers who create GitHub accounts will cease making public commits
within approximately ten years of joining the platform.

The survival fraction is relatively stable from account ages zero to five years
before beginning its decline, which suggests that early dropout is modest and
that the primary attrition occurs in later years. This is consistent with
developers who remain active throughout the early stages of their careers but
gradually reduce open source contribution as professional and other commitments
evolve.

\section{Language-Specific Developer Lifetimes}

We also compute median observed developer lifetimes by language, defined as the
interval between account creation and most recent repository activity.
Table~\ref{tab:lifetime} presents the full results.

\begin{table}[h]
\centering
\begin{tabular}{lrr}
\toprule
Language & Median Lifetime (years) & $n$ \\
\midrule
Elixir       & 12.92 &  32 \\
Clojure      & 12.13 &  52 \\
Emacs Lisp   & 12.03 &  38 \\
Perl         & 11.76 &  54 \\
Nix          & 11.46 &  41 \\
Vim Script   & 10.11 &  60 \\
Lua          & 10.04 & 140 \\
Common Lisp  &  9.50 &  23 \\
Julia        &  9.28 &  50 \\
Swift        &  9.02 & 191 \\
Ruby         &  8.85 & 231 \\
SCSS         &  8.66 & 130 \\
Makefile     &  8.66 &  44 \\
Shell        &  8.62 & 623 \\
Dockerfile   &  8.54 &  61 \\
Go           &  8.35 & 687 \\
Rust         &  8.30 & 490 \\
TeX          &  7.93 &  86 \\
Scala        &  7.69 &  35 \\
HCL          &  7.37 &  52 \\
\bottomrule
\end{tabular}
\caption{Median developer lifetimes by language (languages with $n \geq 23$,
ordered by median lifetime descending).}
\label{tab:lifetime}
\end{table}

These estimates combine two effects: the maturity of the ecosystem and the
retention of developers within that ecosystem. The languages at the top of
Table~\ref{tab:lifetime}---Elixir, Clojure, Emacs Lisp, Perl, and Nix---are
either niche communities with highly committed long-term practitioners or
languages that enjoyed their peak adoption in the early years of GitHub. Older
accounts that remain active have accumulated long measured lifetimes by
construction.

The contrast between median lifetime and activity ratio illuminates different
aspects of ecosystem health. Elixir exhibits a long median lifetime (12.92
years) but does not appear among the highest activity ratios, suggesting that
its developer base is experienced and long-tenured but not rapidly growing.
Conversely, Rust exhibits a relatively modest median lifetime (8.30 years) while
maintaining a very high activity ratio (0.86), consistent with a younger but
highly engaged and growing community.

\section{Discussion}

Even from this partial dataset several structural features of the modern software
ecosystem are legible.

First, newer programming languages such as Rust and TypeScript appear to have
exceptionally high developer retention rates, while ecosystems associated with
earlier waves of GitHub adoption carry larger inactive populations. Second,
developer participation on GitHub follows a characteristic decay pattern with a
half-life of roughly ten years, suggesting that the platform's active developer
population turns over substantially on a decadal timescale. Third, the
distinction between median lifetime and activity ratio reveals that ecosystem
maturity and community vitality are not equivalent: a community can be old and
committed (Elixir, Clojure) or young and growing (Rust, TypeScript) with
importantly different implications for its future trajectory.

The \textsc{None} category deserves further investigation in subsequent work. Its
low activity ratio (0.37) relative to the overall mean (0.60) is consistent with
a hypothesis that inactive developers are less likely to maintain language-tagged
repositories, but it may also reflect structural features of the follow network
itself---for example, spam or bot accounts, which tend to lack repository
content.

\section{Future Work}

The present analysis is based on the first $\sim30{,}000$ accounts processed by
the crawler. The full dataset is expected to contain approximately one million
accounts. Completion of the crawl will enable more precise estimation of:
\begin{itemize}
\item developer survival curves disaggregated by language and account vintage,
\item language ecosystem migration patterns detectable from sequential account activity,
\item the relationship between developer influence (followers) and long-term activity,
\item and formal survival models (e.g., Weibull or log-normal fits) to characterize
      the attrition process quantitatively.
\end{itemize}

Because the dataset is developer-centered rather than repository-centered, it
provides a complementary perspective on the dynamics of open source communities
not available from existing large-scale GitHub datasets.

\section{Conclusion}

Preliminary analysis of a large GitHub follow network suggests that developer
activity on the platform exhibits a characteristic lifetime on the order of a
decade. Activity ratios vary significantly across programming language
ecosystems: Astro, Rust, and TypeScript lead with ratios above 0.85, while
Objective-C, Java, and Ruby trail with ratios at or below 0.49. Median developer
lifetimes are longest in older niche communities (Elixir, Clojure, Emacs Lisp)
and shorter in newer, rapidly growing ecosystems.

Completion of the full dataset will allow a more detailed empirical picture of
developer lifecycles and language ecosystem evolution on GitHub, including formal
survival modeling and analysis of cross-language migration trajectories.

\end{document}
